\chapter{Local Perturbations in Token Space}
\label{ch:rq4}

The previous chapter showed that a local image-space perturbation causes token changes well beyond the corrupted patch. This chapter examines the complementary decoder-side question: when token IDs inside a bounded region of the latent grid are replaced directly, how far does the decoded image change?

The experiment replaces token IDs inside a contiguous rectangular region and
measures the decoded change inside and outside the corresponding image region.
It asks how severe the change is and how far it extends beyond the edited patch.

A discrete image tokenizer represents an image as a spatial grid of codebook
indices. Editing a bounded region of this grid is expected to affect mainly
the corresponding region of the decoded image. Decoder locality matters because
autoregressive image models may make local token errors: if the decoder spreads
such errors widely, a local sequence error can become a broader image artifact.

The replacement strategies test whether nearby codebook vectors decode to
similar image content and whether more distant vectors cause larger changes.
Nearest-neighbour replacement tests a small change in the embedding space;
farthest replacement tests a large one. Orthogonal replacement selects a
vector with cosine similarity near zero, and random replacement provides a
reference independent of distance. We expect these strategies to produce
different levels of change, but their decoded effects must be measured.

\section{Experimental Setup}
\label{sec:rq4_setup}

\subsection{Token-Edit Protocol}
\label{subsec:rq4_token_edit_protocol}

Let $z_0 \in \{1,\ldots,K\}^{h \times w}$ denote the clean token grid produced
for an image, and let $B_{\mathrm{tok}}$ be a contiguous rectangular region of
that grid. For an edit mode $m$, the edited token grid $z^{(m)}$ is
\begin{equation}
    z^{(m)}_{i,j}
    =
    \begin{cases}
        T_m(z_{0,i,j}), & (i,j) \in B_{\mathrm{tok}}, \\
        z_{0,i,j}, & (i,j) \notin B_{\mathrm{tok}},
    \end{cases}
    \label{eq:rq4_token_edit}
\end{equation}
where $T_m$ maps the original token to a replacement token. The experiment
uses four mappings:
\begin{itemize}
    \item \textbf{Closest}: the nearest alternative codebook vector;
    \item \textbf{Orthogonal}: a codebook vector with cosine similarity close
          to zero;
    \item \textbf{Random}: a uniformly sampled codebook vector;
    \item \textbf{Farthest}: the most distant codebook vector.
\end{itemize}

For VQGAN, the
closest, orthogonal, and farthest mappings are restricted to the empirically
active code subset identified in Chapter~\ref{ch:rq1}. VQGAN uses only part of
its codebook on the reference dataset. Uniform-random
replacement is different: it samples from the full codebook for both models.
For VQGAN this means that random replacement may introduce inactive codes, so
the random edit tests both replacement and the use of previously inactive codes.

The clean and edited reconstructions are
\begin{equation}
    x_{\mathrm{clean}} = D(z_0),
    \qquad
    x_{\mathrm{edit}}^{(m)} = D(z^{(m)}),
    \label{eq:rq4_clean_edited_decode}
\end{equation}
where $D$ is the tokenizer decoder. All image changes are measured against
$x_{\mathrm{clean}}$, so ordinary reconstruction error relative to the input
does not enter the perturbation measurement.

The main comparisons use a patch covering $25\%$ of the token grid. A
patch-size sweep also tests target fractions of $10\%$, $50\%$, and $75\%$.
The main analysis uses
ImageNet validation; ImageNet-V2 is used to test whether the observed trends
remain stable under a natural distribution shift.

Images are resized and center-cropped to $256\times256$ using the
model-specific preprocessing described in Appendix~\ref{app:computational_environment}.
Decoder outputs are converted from the models' internal $[-1,1]$ range to
$[0,1]$ before metrics are computed. Patch locations
are reproducible within each run. The two tokenizers did not receive shared
random replacement draws, so comparisons describe aggregate behavior.

Both tokenizers produce a $16 \times 16$ grid. For each image, a square patch
is sampled at a random valid position for the largest target fraction. The
smaller patches are nested within that patch, so comparisons across fractions
do not also change the general image region being edited. The square side lengths are
5, 8, 11, and 14 tokens, corresponding to actual grid fractions of 9.8\%,
25.0\%, 47.3\%, and 76.6\%, respectively.

\subsection{Fidelity Metrics}
\label{subsec:rq4_fidelity_metrics}

Decoder locality alone does not describe the severity of an edit. A low
outside response may coincide with either a negligible local change or severe
damage confined to the patch. Using the definitions and metric references in
Chapter~\ref{ch:datasets_models_metrics}, the experiment measures PSNR, SSIM,
and LPIPS between $x_{\mathrm{clean}}$ and $x_{\mathrm{edit}}^{(m)}$, both over
the full image and within $B_{\mathrm{px}}$. Patch-level metrics quantify edit
severity; the inside/outside response and leakage ratio quantify spatial
propagation. The comparison is made against the clean reconstruction rather
than the original image so that the measurement isolates the decoder response
to the token edit, not the tokenizer's ordinary reconstruction error.

\subsection{Decoder-Locality Metrics}
\label{subsec:rq4_decoder_locality_metrics}

To measure how far decoded changes spread, we compare the average response
inside and outside the edited region. The per-pixel measurement follows the
RMSE construction used in Chapter~\ref{ch:datasets_models_metrics}.
The decoder response is measured relative to the clean reconstruction. For
pixel position $(u,v)$, the channel-averaged root-mean-square change is
\begin{equation}
    \Delta^{(m)}(u,v)
    =
    \sqrt{
        \frac{1}{3}
        \sum_{c=1}^{3}
        \left(
            x_{\mathrm{edit}}^{(m)}(u,v,c)
            -
            x_{\mathrm{clean}}(u,v,c)
        \right)^2
    }.
    \label{eq:rq4_pixel_change}
\end{equation}

Let $B_{\mathrm{px}}$ denote the pixel-space region corresponding to the edited
token patch. Mean inside and outside responses are
\begin{equation}
    \Delta_{\mathrm{in}}^{(m)}
    =
    \frac{1}{|B_{\mathrm{px}}|}
    \sum_{(u,v)\in B_{\mathrm{px}}}
    \Delta^{(m)}(u,v),
    \qquad
    \Delta_{\mathrm{out}}^{(m)}
    =
    \frac{1}{|\overline{B}_{\mathrm{px}}|}
    \sum_{(u,v)\notin B_{\mathrm{px}}}
    \Delta^{(m)}(u,v).
    \label{eq:rq4_inside_outside_change}
\end{equation}
Here $\overline{B}_{\mathrm{px}}$ is the complement of $B_{\mathrm{px}}$
within the image. Its cardinality in the denominator is the number of
\emph{outside} pixels, matching the summation over $(u,v)\notin B_{\mathrm{px}}$.
The decoder leakage ratio is
\begin{equation}
    \lambda_{\mathrm{dec}}^{(m)}
    =
    \frac{\Delta_{\mathrm{out}}^{(m)}}
         {\Delta_{\mathrm{in}}^{(m)} + \varepsilon}.
    \label{eq:rq4_decoder_leakage}
\end{equation}
A small positive $\varepsilon=10^{-8}$ stabilizes the denominator. A low ratio means
that the average pixel change is smaller outside than inside the edited
region. $\Delta_{\mathrm{in}}$ measures local edit strength, while
$\Delta_{\mathrm{out}}$ measures leakage outside the intended image region. The analysis also reports the 99th and 99.9th percentiles of
outside-patch change to capture strong but spatially sparse leakage.

The leakage ratio is calculated per image. Its distribution can become
unstable when
$\Delta_{\mathrm{in}}^{(m)}$ is close to zero: a small outside response divided
by an even smaller inside response can yield a large ratio. The median is
therefore the primary summary of proportional leakage; means and standard
deviations are retained in the archived per-image tables. The absolute inside
and outside responses are reported alongside $\lambda_{\mathrm{dec}}^{(m)}$.

\section{Experimental Results}
\label{sec:rq4_results}

\subsection{Qualitative Decoder Response}
\label{subsec:rq4_qualitative}

Figure~\ref{fig:rq4_qualitative_token_edits} shows one ImageNet validation
example for a 25\% token patch. It shows the clean
reconstruction, the four edited reconstructions, and the corresponding
pixel-change maps for both tokenizers. The nearest-neighbour edit preserves the
fish and surrounding net with only modest local changes. Orthogonal and random
replacement alter the content inside the patch more visibly, while the
farthest edit produces the strongest local transformation for both models.
For VQGAN, the farthest edit becomes an extreme high-contrast artifact,
consistent with the large quantitative response reported below.

\begin{figure}[t]
    \centering
    \includegraphics[width=\textwidth]{ch07_fig_token_edit_qualitative_imagenet_val_patch25.png}
    \caption{Qualitative decoder response to a 25\% token-patch perturbation.
    The cyan rectangle marks the nominal edited region. RMSE heatmaps show
    pixel change relative to each tokenizer's clean reconstruction; the figure
    should be read for the ordering of edit severity and for change outside the
    rectangle, not as a standalone quantitative result.}
    \label{fig:rq4_qualitative_token_edits}
\end{figure}
\subsection{Edit Severity}
\label{subsec:rq4_edit_severity}

Table~\ref{tab:rq4_patch25_severity} and
Figure~\ref{fig:rq4_fidelity_patch25} compare the four replacement strategies
at the 25\% patch level. Closest replacement is the mildest
intervention for both tokenizers. Its patch PSNR is 20.36 dB for LlamaGen and
20.00 dB for VQGAN, while patch LPIPS remains at 0.134 and 0.182. These edits preserve recognizable
content in Figure~\ref{fig:rq4_qualitative_token_edits}; the more distant
replacements visibly change the patch. PSNR measures pixel error, whereas
LPIPS measures differences in learned image features. Moving away
from the nearest code sharply increases perceptual change. Orthogonal
replacement raises patch LPIPS to 0.599 for LlamaGen and 0.614 for VQGAN.
The random edit should be interpreted separately from the three geometric
edits because, for VQGAN, it may sample inactive codes from the full nominal
codebook.

For LlamaGen, orthogonal and random replacement are almost indistinguishable:
their patch LPIPS values are 0.599 and 0.592, and their inside responses are
0.222 and 0.220. VQGAN behaves differently. Random replacement has a lower
mean inside RMSE than orthogonal replacement, 0.200 versus 0.239, but a much
higher patch LPIPS, 0.749 versus 0.614. The disagreement between RMSE and LPIPS
means that average pixel displacement does not account for the perceptual
effect of the VQGAN random edit.

Farthest replacement is the strongest intervention. At the 25\% level,
LlamaGen reaches a patch LPIPS of 0.702 and a patch PSNR of 7.74 dB. VQGAN is
more severely affected under its own farthest-code edit, with patch LPIPS
0.790, patch PSNR 3.74 dB, and an inside response of 0.607. Farthest-code
edits are defined separately in the two codebooks, so this is not a
severity-matched comparison of the decoders.

\begin{table}[t]
    \centering
    \small
    \begin{tabular}{llrrr}
        \textbf{Tokenizer} & \textbf{Edit mode} &
        $\boldsymbol{\Delta_{\mathrm{in}}}$ &
        \textbf{Patch PSNR} & \textbf{Patch LPIPS} \\ \hline
        LlamaGen & Closest & 0.067 & 20.36 & 0.134 \\
        LlamaGen & Orthogonal & 0.222 & 11.45 & 0.599 \\
        LlamaGen & Random & 0.220 & 11.48 & 0.592 \\
        LlamaGen & Farthest & 0.356 & 7.74 & 0.702 \\
        VQGAN & Closest & 0.072 & 20.00 & 0.182 \\
        VQGAN & Orthogonal & 0.239 & 10.66 & 0.614 \\
        VQGAN & Random & 0.200 & 12.83 & 0.749 \\
        VQGAN & Farthest & 0.607 & 3.74 & 0.790 \\
    \end{tabular}
    \caption{Edit severity for a 25\% token patch on ImageNet validation.
    Values are means over 50{,}000 images. The inside RMSE is repeated here to
    relate pixel-response magnitude to patch fidelity; the locality table also
    reports the corresponding outside response.}
    \label{tab:rq4_patch25_severity}
\end{table}
\begin{figure}[t]
    \centering
    \includegraphics[width=\textwidth]{ch07_fig_fidelity_patch25_imagenet_val.png}
    \caption{Patch-level PSNR, SSIM, and LPIPS for the four edit modes at a
    25\% token-patch fraction.}
    \label{fig:rq4_fidelity_patch25}
\end{figure}
\begin{figure}[t]
    \centering
    \includegraphics[width=\textwidth]{ch07_fig_patch_lpips_imagenet_val.png}
    \caption{Patch-level perceptual change as the edited token-patch fraction
    increases.}
    \label{fig:rq4_patch_lpips_sweep}
\end{figure}
Figure~\ref{fig:rq4_patch_lpips_sweep} shows that patch-level perceptual change
generally increases with patch fraction, but the edit-mode ordering remains
stable. Closest replacement stays well below the other interventions across
the sweep. For LlamaGen, orthogonal and random edits continue to track each
other closely, rising from approximately 0.51--0.52 at the 10\% level to
approximately 0.67--0.68 at 75\%. Farthest replacement remains highest,
increasing from 0.630 to 0.782.

VQGAN shows a wider separation between edit modes. Its closest-edit LPIPS
increases from 0.139 to 0.230. Random replacement rises from 0.598 to 0.879,
and overtakes farthest replacement at the 75\% level, where farthest LPIPS is
0.861. The farthest edit still has by far the largest inside RMSE response.
The two metrics rank the strongest VQGAN edit differently.

\subsection{Inside and Outside Decoder Response}
\label{subsec:rq4_inside_outside_results}

Table~\ref{tab:rq4_patch25_locality} separates the response inside the edited
region from the response outside it. At the 25\% patch level, all edit modes
produce a substantially larger mean change inside the patch. For LlamaGen,
outside response ranges from 0.013 for closest replacement to 0.050 for
farthest replacement. For VQGAN, it ranges from 0.019 to 0.090. The decoder
response is spatially concentrated, but not confined to the edited region.

The heatmaps in Figure~\ref{fig:rq4_qualitative_token_edits} show the same
pattern: the response does not stop exactly at the cyan boundary. Most change
remains within or near the edited region, but lower-amplitude changes extend
into the surrounding image, including under closest replacement.

Even nearest-neighbour replacement changes pixels outside the patch, as the
heatmaps in Figure~\ref{fig:rq4_qualitative_token_edits} also show. Farthest
replacement causes the greatest absolute change both inside and outside.
Its small ratio therefore reflects severe inside damage and should not be
read as evidence that the replacement is harmless outside the patch.

The outside response also contains localized high-amplitude changes that are
not visible from its mean alone. At the 25\% patch level, VQGAN farthest
replacement has a mean outside response of 0.090, but its mean outside 99th and
99.9th percentiles reach 0.673 and 0.854. For LlamaGen, the corresponding
values are 0.050, 0.328, and 0.554. Closest replacement has much smaller
outside tails, with 99th-percentile values of 0.105 for LlamaGen and 0.148 for
VQGAN. These tail statistics matter because mean outside response can hide
localized artifacts: strong changes can occur at a small subset of outside
pixels even when the average outside response remains low.

\begin{table}[t]
    \centering
    \small
    \begin{tabular}{llrrr}
        \textbf{Tokenizer} & \textbf{Edit mode} &
        $\boldsymbol{\Delta_{\mathrm{in}}}$ &
        $\boldsymbol{\Delta_{\mathrm{out}}}$ &
        $\boldsymbol{\widetilde{\lambda}_{\mathrm{dec}}}$ \\ \hline
        LlamaGen & Closest & 0.067 & 0.013 & 0.193 \\
        LlamaGen & Orthogonal & 0.222 & 0.039 & 0.176 \\
        LlamaGen & Random & 0.220 & 0.038 & 0.175 \\
        LlamaGen & Farthest & 0.356 & 0.050 & 0.141 \\
        VQGAN & Closest & 0.072 & 0.019 & 0.263 \\
        VQGAN & Orthogonal & 0.239 & 0.046 & 0.187 \\
        VQGAN & Random & 0.200 & 0.050 & 0.249 \\
        VQGAN & Farthest & 0.607 & 0.090 & 0.148 \\
    \end{tabular}
    \caption{Mean inside and outside response and median decoder leakage ratio
    for a 25\% token patch on ImageNet validation.}
    \label{tab:rq4_patch25_locality}
\end{table}
\begin{figure}[t]
    \centering
    \includegraphics[width=\textwidth]{ch07_fig_decoder_locality_imagenet_val.png}
    \caption{Mean image-space response inside and outside the edited token
    patch as the patch fraction increases.}
    \label{fig:rq4_decoder_locality}
\end{figure}
\begin{figure}[t]
    \centering
    \includegraphics[width=\textwidth]{ch07_fig_leakage_ratio_imagenet_val.png}
    \caption{Median decoder leakage ratio as a function of token-patch
    fraction and edit mode.}
    \label{fig:rq4_leakage_ratio}
\end{figure}
Figure~\ref{fig:rq4_decoder_locality} shows how these quantities change with
patch size. The mean inside response is comparatively stable for closest,
orthogonal, and farthest replacement. For example, VQGAN's farthest-edit
inside response changes only from 0.592 at the 10\% level to 0.613 at 75\%.
LlamaGen's corresponding response rises from 0.329 to 0.410. By contrast, the
outside response increases with patch size. For VQGAN farthest
replacement, the outside response rises from 0.053 to 0.213; for LlamaGen it
rises from 0.030 to 0.120.

The same growth appears in Figure~\ref{fig:rq4_leakage_ratio}. At the 75\%
level, median leakage ratios range from 0.299 to 0.418 for LlamaGen and from
0.352 to 0.531 for VQGAN. The closest and random VQGAN edits have the highest
proportional leakage at large patch sizes. The plots use median per-image
ratios throughout; near-zero inside responses can distort the mean.

\subsection{Robustness under Dataset Shift}
\label{subsec:rq4_dataset_shift}

We test ImageNet-V2 to check whether the measured decoder locality persists
beyond the in-distribution reference images. Patch-LPIPS differences from the
ImageNet validation baseline are small across both models, all four edit
modes, and all patch fractions: the largest absolute difference is 0.0068.
This is too small to support an interpretation of the individual signed
differences.

The other metrics are similarly stable. Across the full comparison, the
largest absolute change in mean outside response is 0.0016. At the 25\% patch
level, leakage-ratio differences range from $-0.0030$ to $-0.0080$, and
patch-PSNR differences remain within 0.17 dB. These differences are small
relative to those among edit modes and tokenizers. ImageNet-V2 preserves the
ordering and approximate magnitude found on ImageNet validation.
\section{Discussion}
\label{sec:rq4_discussion}

A local token-space perturbation produces a spatially concentrated decoder
response, but the response is not strictly local. Mean change is always stronger
inside the corresponding image region than outside it, and the qualitative
heatmaps preserve a clear spatial relationship to the edited patch.
Outside response is measurable even for nearest-neighbour replacement and grows as
the edited region becomes larger. Farthest-code replacement causes large
absolute changes both inside and outside the patch.

Orthogonal and random replacements are nearly equivalent for LlamaGen.
For VQGAN, random replacement is more disruptive under LPIPS than the
orthogonal edit despite a smaller average pixel response.

VQGAN random replacement can introduce codes absent from its clean ImageNet
assignments. This affects the interpretation of its scores and qualitative
examples: the experiment does not isolate the effect of distance from the
effect of using inactive codes.

The same within-checkpoint edit patterns persist across patch fractions.
At the 25\% level, the VQGAN farthest-code edit has an inside response of
0.607, compared with 0.356 for LlamaGen's separately defined farthest-code
edit; its patch PSNR is approximately 4 dB lower. These edits are not
severity-matched across codebooks, so the difference does not isolate decoder
sensitivity. Neither decoder confines the effect completely.

These decoder-side findings contrast with the encoder-side results in
Chapter~\ref{ch:rq3}. A local image-space perturbation caused widespread token
identity changes outside the patch, with outside flip rates approaching one
half under stronger perturbations. Direct token edits produce a more localized
image response: outside RMSE remains much smaller than inside RMSE.
The two experiments use different response variables, so their magnitudes are
not directly comparable. They nevertheless locate two separate effects: local
image perturbations can alter token identities far outside the patch, while
direct token edits produce a decoded response that is concentrated near the
edited region but still extends beyond it.

ImageNet-V2 preserves the same ordering and approximate response magnitude.
Decoder locality therefore persists under this mild shift to natural images
from different sources. Testing whether it persists under stronger style and
domain changes remains future work.
